% Part: turing-machines
% Chapter: machines-computations
% Section: turing-machines

\documentclass[../../../include/open-logic-section]{subfiles}

\begin{document}

\olfileid[hi]{tur}{mac}{tur}
\olsection{ट्यूरिंग मशीनें}

\begin{explain}
ट्यूरिंग मशीन किन वस्तुओं से बनती है, इसकी आकारिक परिभाषा अमूर्त दिखती
है, किंतु वास्तव में सरल है: वह ट्यूरिंग मशीन की कार्यप्रणाली निर्दिष्ट
करने के लिए आवश्यक सारी सूचना को केवल एक गणितीय संरचना में समेटती है।
इसमें यह सम्मिलित है कि (1) मशीन किन अवस्थाओं में हो सकती है, (2) पट्टी
पर किन प्रतीकों की अनुमति है, (3) मशीन को किस अवस्था में आरंभ होना चाहिए,
और (4) मशीन का निर्देश-समुच्चय क्या है।
\end{explain}

\begin{defn}[ट्यूरिंग मशीन]
एक \emph{ट्यूरिंग मशीन} $M$ एक टपल $\langle Q, \Sigma, q_0,
\delta\rangle$ है, जिसमें होते हैं:
\begin{enumerate}
\item \emph{अवस्थाओं} का एक परिमित समुच्चय~$Q$,
\item एक परिमित \emph{वर्णमाला} $\Sigma$, जिसमें $\TMendtape$ और
  $\TMblank$ सम्मिलित हैं,
\item एक \emph{आरंभिक अवस्था}~$q_0 \in Q$,
\item एक परिमित \emph{निर्देश-समुच्चय}~$\delta\colon Q \times \Sigma
  \pto Q \times \Sigma \times \{\TMleft, \TMright, \TMstay\}$.
\end{enumerate}
आंशिक फलन~$\delta$ को $M$ का \emph{संक्रमण फलन} भी कहते हैं।
\end{defn}

\begin{explain}
हम मानते हैं कि पट्टी केवल एक दिशा में अनंत है। इसलिए पट्टी के बाएँ अंत
के चिह्नक के रूप में एक विशेष प्रतीक~$\TMendtape$ निर्दिष्ट करना उपयोगी
है। इससे ट्यूरिंग मशीन के प्रोग्राम के लिए यह बताना सरल हो जाता है कि वह
कब पट्टी से बाहर निकलने के ``संकट'' में है। हम मान सकते थे कि इस प्रतीक
के ऊपर कभी नहीं लिखा जाता; अर्थात्
$\delta(q,\TMendtape) = \tuple{q', \TMendtape, x}$, यदि
$\delta(q,\TMendtape)$ परिभाषित हो। कुछ
पाठ्यपुस्तकें ऐसा करती हैं; हम नहीं करते। अपनी ट्यूरिंग मशीन बनाते समय आप
केवल इतना ध्यान रख सकते हैं कि वह~$\TMendtape$ के ऊपर कभी न लिखे। इसके
अतिरिक्त, कुछ स्थितियों में इसके ऊपर लिखने की अनुमति सुविधाजनक लचीलापन
देती है।
\end{explain}

\begin{ex}
\emph{सम मशीन:} सम मशीन आकारिक रूप से चतुष्टय
$\tuple{Q, \Sigma, q_0, \delta}$ है, जहाँ
\begin{align*}
Q & = \{ q_0, q_1 \} \\
\Sigma & = \{ \TMendtape, \TMblank, \TMstroke \}, \\
\delta(q_0, \TMstroke) & = \tuple{q_1, \TMstroke, \TMright},\\
\delta(q_1, \TMstroke) & = \tuple{q_0, \TMstroke, \TMright},\\
\delta(q_1, \TMblank)  & = \tuple{q_1, \TMblank, \TMright}.
\end{align*}
\end{ex}

\end{document}
